% Generado por generar_capitulos.py a partir de ../.tex8; no editar a mano
\chapter{Marco Teórico}
\label{sec:marco_teorico}
% <- 02
\section{Memoria Asociativa Entrópica}
\label{sec:eam}

La Memoria Asociativa Entrópica (EAM) \cite{pineda2021entropic} es un modelo de memoria
que opera sobre funciones discretas.
Una \textbf{pista}, lo que se presenta a la memoria para registrarlo o para preguntar por
él, es una función discreta $f : \{1, \ldots, n\} \to \{0, \ldots, m-1\}$,
representada como un vector $\mathbf{a} \in \{0, \ldots, m-1\}^n$.
La dimensión $n$ es el número de \emph{características} y $m$ es el número de valores posibles
por característica (el alfabeto).

En la formulación original de 2021 el registro es \emph{booleano}; la memoria almacena
una función característica $R : A \times V \to \{0, 1\}$, sobre características ($A$) y
valores ($V$), y el registro se realiza por disyunción lógica.
La variante que este trabajo utiliza es la \textbf{EAM ponderada} (W-EAM), introducida
en el aprendizaje fonético \cite{pineda2022weighted} y empleada desde entonces en la
línea de trabajo del grupo, de la imaginería a las imágenes naturales de
CIFAR-10 \cite{pineda2023imagery,hernandez2026cifar}.
En ella, la memoria se representa como una relación ponderada
$\mathbf{R} \in \mathbb{N}^{n \times m}$, donde $\mathbf{R}[i, j]$ cuenta cuántas veces la característica $i$ ha tomado el valor
$j$ a lo largo de todos los registros.
Las definiciones que siguen corresponden a la variante ponderada, no a la booleana de
2021.

\subsection*{Operación de registro ($\lambda$)}

Dada una pista $\mathbf{a}$, se construye el patrón unitario:
\begin{equation}
    r(\mathbf{a})[i, j] = \begin{cases} 1 & \text{si } j = a_i \\ 0 & \text{en otro caso} \end{cases}
    \label{eq:patron-unitario}
\end{equation}
y se acumula en la memoria:
\begin{equation}
    \mathbf{R} \leftarrow \mathbf{R} + r(\mathbf{a})
    \label{eq:registro}
\end{equation}
La acumulación está acotada superiormente, tal como definen la W-EAM original
\cite{pineda2022weighted} y la fuente heteroasociativa \cite{morales2025missing}, cuyos
autores describen esta operación como una forma de la regla de aprendizaje de Hebb.

La memoria acumula una \textbf{distribución de experiencias}.
Cada vector $\mathbf{R}[i, \cdot]$ (la \emph{columna} de la característica $i$, como la
llaman las fuentes) se convierte progresivamente en una distribución de
frecuencias sobre los valores posibles de la característica $i$.
De esta superposición estadística van saliendo abstracciones y prototipos, sin que se
impongan desde el exterior ni se calculen por separado.

\subsection*{Operación de reconocimiento ($\eta$)}

Se dice que la pista $\mathbf{a}$ está \emph{contenida} en $\mathbf{R}$ si, para toda
característica $i$, la celda correspondiente tiene masa positiva: $\mathbf{R}[i, a_i] > 0$.
La versión con tolerancia $\xi$ admite a lo más $\xi$ fallas:
\begin{equation}
    \eta(\mathbf{a}, \mathbf{R}, \xi) = \mathbf{1}\!\left[\,\bigl|\{i : \mathbf{R}[i, a_i] = 0\}\bigr| \leq \xi\,\right]
    \label{eq:reconocimiento}
\end{equation}
donde $\mathbf{1}[\cdot]$ vale 1 si la condición se cumple y 0 si no.
Es una \emph{implicación material relajada}: la pista se acepta si cada una de sus
celdas marcadas encuentra soporte en la memoria, tolerando hasta $\xi$ columnas sin él.

Durante el reconocimiento se computan además los \textbf{pesos por dimensión}, que este
trabajo escribe como la función $\omega$:
\begin{equation}
    \omega(\mathbf{a}, \mathbf{R}) = (w_1, \ldots, w_n), \qquad w_i = \mathbf{R}[i,\, a_i]
    \label{eq:pesos}
\end{equation}
Estos pesos reflejan cuán frecuentemente la característica $i$ ha tenido el valor $a_i$
en los patrones registrados.
En las definiciones formales de las fuentes, $\eta$ devuelve solo aceptar o rechazar, y
los $w_i$ son cantidades intermedias de esa prueba; la implementación de referencia de
Pineda y Morales los devuelve\footnote{Función \texttt{recog\_weights}.}, y este
trabajo los usa para ponderar la proyección de la memoria heteroasociativa (Sección~\ref{sec:hetero4d}).

El parámetro $\iota$ controla la poda.
Para cada característica $i$ se toma el peso medio de sus celdas no nulas,
$\bar{r}_i = \sum_j \mathbf{R}[i,j] \,/\, |\{j : \mathbf{R}[i,j] > 0\}|$, y toda celda
con $\mathbf{R}[i,j] < \iota\,\bar{r}_i$ se apaga al reconocer y al recuperar.
La relación acumulada $\mathbf{R}$ no cambia, porque la poda quita los valores poco
frecuentes solo de la \emph{relación operativa} con la que la memoria responde.

El parámetro $\kappa$ fija un umbral de reconocimiento relativo: la memoria
homoasociativa acepta la pista solo si, además de contenerla, el peso medio
$\bar{w} = \frac{1}{n}\sum_i w_i$ alcanza $\kappa \cdot \bar{\mu}$, donde $\bar{\mu}$ es
la media global de la memoria.
Este sistema usa de esta memoria los pesos, pero no esa decisión
(Sección~\ref{sec:hetero4d}); el $\kappa$ que barre la
Sección~\ref{sec:param-iota-kappa} actúa sobre la proyección heteroasociativa.

\subsection*{Operación de recuperación ($\beta$)}

La recuperación de un patrón a partir de una pista $\mathbf{a}$ procede columna a columna.
Para cada característica $i$, se construye una distribución de probabilidad sobre $\{0,\ldots,m-1\}$
modulada por una gaussiana de desviación $\sigma \cdot m$ centrada en el valor $a_i$
(cuando se dispone de la pista):
\begin{equation}
    p_i(j) \propto \mathbf{R}[i, j] \cdot \exp\!\left(-\frac{(j - a_i)^2}{2(\sigma m)^2}\right)
    \label{eq:distribucion-recuperacion}
\end{equation}
El valor recuperado $\hat{a}_i$ se obtiene \textbf{muestreando} esta distribución:
$\hat{a}_i \sim p_i(\cdot)$.
Por eso dos ejecuciones de $\beta$ sobre la misma pista y la misma memoria pueden dar
resultados distintos; esa variación es parte de cómo recuerda la memoria:
\emph{recordar es muestrear}, no leer un valor fijo.
El parámetro $\sigma$ controla la dispersión del muestreo.
Con $\sigma = 0$ se concentra la masa en la pista exacta, valores mayores permiten
explorar vecindades del espacio de representación; valores altos favorecen
recuperaciones imaginativas a costa de falsos recuerdos \cite{pineda2023imagery}.
Esta modulación gaussiana es la recuperación de la W-EAM \cite{pineda2022weighted}
(la EAM de 2021 usaba una distribución triangular \cite{pineda2021entropic}) y es la
que implementa el código de referencia que este sistema utiliza.

La \textbf{entropía} de la memoria se calcula como la media de las entropías de Shannon
por columna:
\begin{equation}
    H(\mathbf{R}) = \frac{1}{n} \sum_{i=1}^{n} \left(-\sum_{j=0}^{m-1} \hat{p}_{ij} \log_2 \hat{p}_{ij}\right)
    \label{eq:entropia}
\end{equation}
donde $\hat{p}_{ij} = \mathbf{R}[i,j] / \sum_l \mathbf{R}[i,l]$.
La entropía resume el estado de la memoria.
Una memoria con entropía alta ha acumulado más variedad; una con entropía baja está
concentrada en pocos valores.
La entropía baja \emph{no} es un defecto por sí misma: en las
fuentes, una función totalmente determinada tiene entropía cero y es una
memoria bien formada, solo que rígida y reproductiva
\cite{pineda2021entropic,morales2022manuscript}.
Si la entropía baja es deseable o no depende de la tarea.
Un directorio que debe señalar sin ambigüedad a un especialista necesita poca
incertidumbre sobre el módulo dada la pista (Sección~\ref{sec:tmt}), y una memoria de
contenido que debe generalizar, un rango medio de entropía.
La evidencia empírica de las fuentes ubica el óptimo de reconocimiento y recuperación
constructiva en un rango \emph{intermedio} de entropía.
Con entropía baja, el reconocimiento comete pocas aceptaciones falsas y muchas
omisiones; con entropía alta ocurre lo contrario \cite{pineda2023imagery}.

\section{Memoria Heteroasociativa Entrópica}
\label{sec:hetero4d}

La Memoria Heteroasociativa Entrópica (EHAM) de Morales y
Pineda~\cite{morales2025missing}
extiende la EAM al almacenamiento de pares de objetos posiblemente pertenecientes a
dominios o modalidades distintos.
En lugar de una relación bidimensional, la EHAM almacena las asociaciones en una relación
ponderada de cuatro dimensiones, de modo que una pista en un dominio puede proyectar una
relación en el dominio asociado.
Formalmente, relaciona un dominio izquierdo con $n$ características y alfabeto de tamaño $m$,
y un dominio derecho con $p$ características y alfabeto de tamaño $q$.
En el artículo la relación es $h \subseteq A \times B \times V \times Z$, con pesos
acotados; la implementación de referencia la almacena en un tensor
$\mathbf{R} \in \mathbb{N}^{n \times p \times (m+1) \times (q+1)}$,
donde los tamaños $(m+1)$ y $(q+1)$ reservan la última posición de cada alfabeto para
un \emph{valor indefinido}, que representa características sin valor (funciones
parciales); esas celdas forman el \emph{borde} del tensor.
Las operaciones $\lambda$, $\eta$ y $H$ se generalizan de forma directa: el registro es
la misma suma acotada sobre el tensor 4D, el reconocimiento es la misma implicación
material relajada, con la poda $\iota$ y el umbral $\kappa$, y la entropía es el
promedio, sobre los pares de características $(i, j)$, de la entropía de Shannon del plano
$\mathbf{R}[i, j, \cdot, \cdot]$ \cite{morales2025missing}.

La recuperación heteroasociativa introduce además el \textbf{problema de la pista faltante}
(\emph{missing cue problem}), contribución central del artículo publicado
\cite{morales2025missing}.
Cuando la pista fuente selecciona un plano de memoria en el dominio destino pero no
existe información adicional de ese dominio, el plano queda en gran parte indeterminado y no hay con qué completar la recuperación.
Los autores proponen los métodos \emph{random}, \emph{sample and test} y
\emph{search and test} para ese caso.
La proyección ponderada que se describe a continuación produce ese plano a partir de la
pista fuente y sus pesos; para elegir un objeto del plano, este sistema usa el tercer
método, \emph{search and test}, que es la recuperación por defecto de la implementación
de referencia\footnote{Función \texttt{recall\_with\_sampling\_n\_search}.}.
El presente sistema usa la EHAM para la memoria de contenido de cada módulo, que asocia
texto e imagen en los dos sentidos, y para sus directorios, que asocian una pista (una
palabra o una imagen) con la identidad del módulo que ganó la consulta.

\subsection*{Registro heteroasociativo}

Dada una pista izquierda $\mathbf{a} \in \{0,\ldots,m-1\}^n$ y una pista derecha
$\mathbf{b} \in \{0,\ldots,q-1\}^p$, el patrón unitario conjunto es:
\begin{equation}
    r(\mathbf{a}, \mathbf{b})[i, j, a_i, b_j] = 1
    \label{eq:registro-hetero}
\end{equation}
y se acumula de forma idéntica a la ecuación~\eqref{eq:registro}.

\subsection*{Las memorias de dominio como lecturas de la heteroasociativa}

La memoria heteroasociativa modela la \emph{distribución conjunta} de dos dominios, y
la de cada dominio se lee de ella.
Cada registro de la ecuación~\eqref{eq:registro-hetero} suma una unidad en
$(i, j, a_i, b_j)$ para \emph{todas} las coordenadas $j$ del otro dominio.
Por eso, fijada cualquier $j$, sumar sobre los valores $l$ del otro dominio cuenta
cuántas veces se registró cada valor $k$ en cada característica $i$, sea cual sea la $j$
elegida:
\begin{equation}
    M^{L}[i, k] = \sum_{l=0}^{q} \mathbf{R}[i, j, k, l],
    \qquad
    M^{R}[j, l] = \sum_{k=0}^{m} \mathbf{R}[i, j, k, l].
    \label{eq:marginal}
\end{equation}
El índice $l = q$ (o $k = m$) es el valor indefinido, el borde.
Lo que llega sin pareja (una imagen sin palabra, una palabra sin imagen) se registra ahí
y entra en la suma.
Las celdas del borde se inicializan en 1, y la lectura descuenta ese 1.
Estas lecturas, $M^{L}$ y $M^{R}$, son las relaciones de memorias homoasociativas
llenadas con lo mismo que recibió cada lado; la prueba $\eta$, los pesos de la
ecuación~\eqref{eq:pesos} y la recuperación dentro de un dominio se calculan sobre ellas.

Por eso cada módulo de EAM-TMS tiene una sola memoria de contenido, la
heteroasociativa, y dos directorios heteroasociativos (uno de texto y uno visual).
Un módulo que manejara un solo dominio guardaría su contenido en una memoria
homoasociativa.
En la implementación de referencia de Pineda y
Morales\footnote{\url{https://github.com/eam-experiments/hetero}, archivo \texttt{eam.py}.},
las memorias homoasociativas se llenan aparte y solo entregan los pesos de la prueba de
reconocimiento de pares; la recuperación usa pesos uniformes.
La Sección~\ref{sec:una-memoria} muestra que leer esos pesos de la heteroasociativa
produce las mismas decisiones.

\subsection*{Proyección AND-ponderada}

Sea $\mathbf{a}$ una pista izquierda con pesos por dimensión $\mathbf{w}$, calculados
con la ecuación~\eqref{eq:pesos} sobre la lectura $M^{L}$ de la propia memoria
(ecuación~\eqref{eq:marginal}).
Se normalizan los pesos:
\begin{equation}
    w'_i = \frac{n \cdot w_i}{\displaystyle\sum_{k=1}^{n} w_k}
    \label{eq:normalizacion-pesos}
\end{equation}
de modo que $\frac{1}{n}\sum_i w'_i = 1$ (media unitaria).
La \textbf{proyección} $\Pi(\mathbf{a})$ sobre el dominio derecho se define por:
\begin{equation}
    \Pi[j, l] = \begin{cases}
        0 & \text{si } \exists\, i : \mathbf{R}[i, j, a_i, l] = 0 \\[4pt]
        \displaystyle\sum_{i=1}^{n} w'_i\, \mathbf{R}[i, j, a_i, l] & \text{en otro caso}
    \end{cases}
    \label{eq:proyeccion-AND}
\end{equation}
La condición de anulación es la \textbf{conjunción AND}: una celda del resultado se
anula si \emph{cualquiera} de las $n$ evidencias aporta cero.
Esto implementa una \emph{contención} estricta, porque la proyección solo produce masa
donde todas las características de la pista tienen soporte simultáneo en la memoria.
Los pesos $w'_i$ ponderan la suma por la frecuencia con que cada característica de la
pista se ha registrado en la memoria, de modo que las características muy frecuentes
pesan más.

El nombre «proyección AND-ponderada» es de este trabajo; el mecanismo no.
Es la operación de proyección de la implementación de referencia de Pineda y
Morales\footnote{Función \texttt{project()}.}, que este sistema incorpora sin
modificaciones.
La anulación conjuntiva y la suma ponderada por pesos normalizados son líneas de esa
función, aplicadas sobre la relación ya podada por $\iota$.
La definición escrita de la recuperación $\beta$ en \cite{morales2025missing} suma las
rebanadas de todas las características sin anular celdas (una unión ponderada); la
anulación conjuntiva es de la implementación, y es la que usan aquí tanto la
recuperación como las lecturas de directorio.

La recuperación heteroasociativa $\beta$ muestrea $\Pi$ columna a columna varias veces,
conserva la muestra que, proyectada de regreso, más se parece a la pista, y la mejora
con una búsqueda local en su vecindad; es el método \emph{search and test} de la pista
faltante.

La memoria heteroasociativa de referencia define el umbral $\kappa$ solo para un par de
pistas.
Con una sola pista, este trabajo lo aplica a la proyección y compara la activación media
por celda de $\Pi$, por coordenada de la pista, con $\kappa \cdot \bar{\mu}$ de la
memoria de contenido.

\section{Memoria transactiva}
\label{sec:tmt}

Wegner~\cite{wegner1987transactive} define un sistema de memoria transactiva como un
conjunto de sistemas de memoria individuales en combinación con la comunicación que
ocurre entre los individuos (p.~186).
La primera mitad de la definición es la memoria de cada miembro; la segunda, lo que
hace de esas memorias un sistema.
En su modelo posterior, Wegner~\cite{wegner1995computer} compara el grupo con una red
de computadoras en la que cada miembro guarda parte de lo que el grupo sabe y lleva un
\emph{directorio} de quién sabe qué, separado del contenido.
En la red distingue tres procesos: la actualización del directorio (\emph{directory
updating}), en la que los miembros aprenden quién sabe qué en el grupo; la asignación
de información (\emph{information allocation}), que asigna cada ítem de memoria a un
miembro, y la coordinación de la recuperación (\emph{retrieval coordination}), que
planea cómo encontrar cada ítem.
La Tabla~\ref{tab:wegner-sistema} separa lo que este trabajo toma de cada fuente de lo
que decide EAM-TMS.

\begin{table}[H]
\centering
\caption[Lo que EAM-TMS toma de Wegner y lo que decide el sistema]{Lo que EAM-TMS toma de Wegner y lo que decide el sistema.
«---»: el elemento no se toma de esa fuente.}
\label{tab:wegner-sistema}
\small
\setlength{\tabcolsep}{4pt}
\begin{tabular}{@{}p{2.6cm}p{3.1cm}p{3.6cm}p{5.6cm}@{}}
\toprule
& \textbf{Wegner (1987)} & \textbf{Wegner (1995)} & \textbf{Decisión de EAM-TMS} \\
\midrule
Sistema & memorias individuales y la comunicación entre individuos (p.~186) &
  memorias individuales enlazadas como computadoras en red &
  ocho módulos con una memoria de contenido y dos directorios cada uno; la
  comunicación es la consulta que un módulo reenvía a otro \\
Directorio & --- & lo que cada miembro sabe de quién sabe qué &
  memoria heteroasociativa de la pista a la identidad del módulo; perspectival en el
  sistema y en la comparación con una sola EHAM, compartido en los demás estudios de
  texto (Sección~\ref{sec:organizacion}) \\
Actualización del directorio & --- & aprender quién sabe qué en el grupo &
  registrar (pista, ganador) tras cada consulta resuelta; solo por interacción \\
Asignación de información & --- & asignar cada ítem de memoria a un miembro &
  no se modela; la etiqueta de clase decide qué módulo guarda cada imagen y cada palabra \\
Coordinación de la recuperación & --- & planear cómo encontrar cada ítem &
  el módulo de entrada suma su directorio y los de sus conocidos (lectura normalizada
  por conteo) y reenvía al de mayor puntuación \\
Mediador y fases & --- & --- &
  el TME media en la fase temprana; el paso a la fase madura lo fija el protocolo \\
Rechazo & --- & --- &
  si nadie tiene soporte para la consulta, el sistema no responde \\
\bottomrule
\end{tabular}
\end{table}

En este sistema, el directorio se actualiza solo por interacción, y lo anotan los
módulos que participaron en cada consulta resuelta (Sección~\ref{sec:componentes}).
En los grupos de personas no es el único camino, porque los miembros también se
reparten las categorías que recordará cada uno según estereotipos de
género~\cite{hollingshead2003gender}, y los grupos que reciben información sobre las
habilidades de sus miembros, sin haber entrenado juntos, rinden como los que entrenaron
juntos~\cite{moreland2000exploring}.
La asignación de información no se modela; la hace la etiqueta de clase.
La especialización de cada módulo es por tanto asignada, como en la condición de
estructura asignada de Wegner, Erber y Raymond~\cite{wegner1991transactive}, y lo que
se forma con el uso es el directorio.
En la coordinación de la recuperación, cuando nadie tiene soporte para la consulta, el
grupo la rechaza; en este sistema, el rechazo es una respuesta válida.

Las dos fases son también una decisión de EAM-TMS.
En la \textbf{fase temprana} el grupo depende de un mediador, porque los miembros no
conocen aún las especialidades de los demás y las consultas se difunden a todos.
En la \textbf{fase madura}, cada módulo enruta con su directorio y los de sus
conocidos, y el mediador sale del circuito.
El paso de una a otra lo fija el protocolo experimental, tras un número dado de
interacciones; un criterio interno de madurez queda como trabajo futuro.

La \textbf{entropía del directorio} resume cómo se reparten los registros entre
módulos (Sección~\ref{sec:operaciones}); no es la entropía de una memoria
(ecuación~\eqref{eq:entropia}).
En el registro completo del grupo (el del mediador, o el directorio compartido de la
calibración), el máximo, $\log_2 8 = 3$ bits con ocho módulos, indica que los registros
están repartidos de forma pareja, y una entropía baja indica colapso hacia un solo
módulo o formación incompleta.
Mide balance, no especialización.
Con un banco balanceado entre clases, un directorio
que anotara ganadores al azar también quedaría cerca del máximo.
Que el grupo esté especializado lo dice la relación entre el dominio de cada consulta y
el módulo que la ganó, que se mide con la información mutua entre ambos
(Sección~\ref{sec:operaciones}); es lo más cercano, en este sistema, a la
especialización que se mide en los grupos de personas~\cite{austin2003transactive}.
En el directorio de cada módulo el máximo no es el estado deseable, porque cada uno
registra lo que presenció y su directorio se carga hacia sí mismo y queda por debajo del
máximo aunque el grupo esté balanceado.

\section{Trabajo relacionado}
\label{sec:trabajo-relacionado}

Repartir el conocimiento entre especialistas y elegir cuál responde tiene varias
tradiciones.

\paragraph{Mezcla de expertos.}
En una mezcla de expertos (\textit{mixture of experts}), una red de compuerta decide qué red
especialista atiende cada caso, y cada especialista termina por cubrir una región
distinta de los datos~\cite{jacobs1991adaptive}, con variantes en árbol entrenadas con el algoritmo
EM~\cite{jordan1994hierarchical} y, a gran escala, compuertas que eligen para cada
ejemplo pocos expertos~\cite{shazeer2017outrageously} o uno solo~\cite{fedus2022switch}.
EAM-TMS comparte la estructura pero no el aprendizaje.
Sus especialistas son memorias
llenadas por clase, y su directorio se forma registrando quién resolvió cada consulta,
sin gradientes.
Las mezclas dispersas agregan una pérdida auxiliar para repartir la carga entre
expertos~\cite{fedus2022switch}; la lectura normalizada por conteo
(Sección~\ref{sec:operaciones}) reencuentra ese problema, el del módulo que gana por
tener más registros, y lo resuelve dividiendo entre el conteo.

\paragraph{Memorias asociativas.}
La EAM pertenece a la familia de la memoria asociativa no holográfica de Willshaw y
colaboradores~\cite{willshaw1969nonholographic}, las memorias asociativas bidireccionales
de Kosko~\cite{kosko1988bidirectional}, la red de Hopfield, que recupera un recuerdo
completo a partir de una parte suficiente~\cite{hopfield1982neural}, la memoria
distribuida dispersa~\cite{kanerva1988sparse}, las memorias densas, de mayor
capacidad~\cite{krotov2016dense}, y las redes de Hopfield modernas, de estados
continuos~\cite{ramsauer2021hopfield}.
Se distingue por guardar distribuciones de frecuencia por característica, reconocer
por contención y recordar muestreando (Sección~\ref{sec:eam}).
Cuando una sola memoria guarda muchos patrones aparecen estados metaestables que los
mezclan~\cite{amit1985spin,ramsauer2021hopfield}, y una red conexionista que aprende tareas una tras otra puede perder lo que ya sabía
(interferencia catastrófica)~\cite{mccloskey1989catastrophic}.
La ventaja de partir el contenido entre especialistas (Sección~\ref{sec:eham-unica})
responde al mismo principio.
Con un dominio por memoria, al recordar se mezclan menos
patrones de otros dominios.

\paragraph{Memoria transactiva computacional.}
La memoria transactiva ya se ha simulado.
En un modelo computacional, saber qué saben los demás acelera más el trabajo de los
grupos grandes y de los que enfrentan tareas o conocimiento cambiantes, mientras que en la
calidad de las decisiones ganan más los grupos pequeños~\cite{ren2006contingent}.
Otros modelan la coevolución de las
redes de comunicación y de conocimiento de un grupo~\cite{palazzolo2006coevolution}.
EAM-TMS no simula un grupo para estudiar su desempeño; lo construye con memorias
asociativas que resuelven consultas de texto e imagen.
Este puente entre la teoría de Wegner y las memorias asociativas entrópicas es una
contribución de este trabajo.
Los artículos de la línea EAM/EHAM no contienen nociones de directorio, de especialistas que cooperan ni de memoria transactiva.

\paragraph{Sistemas multiagente.}
En la red de contratos (\textit{contract net}), las tareas se reparten por negociación.
Un gestor anuncia la tarea, los nodos ofertan y el gestor adjudica~\cite{smith1980contract},
como en la fase temprana, donde el TME difunde la consulta y gana el módulo con más
soporte.
Los intermediarios (\textit{middle-agents}) procesan anuncios de lo que ofrece
cada proveedor, solicitudes o ambas cosas, y median entre quien pide y quien
provee~\cite{decker1997middle}; el TME se les parece, pero aquí nadie anuncia lo que
sabe, y el registro se llena observando quién resolvió cada consulta.

\paragraph{Enrutamiento distribuido y búsqueda federada.}
En las redes entre pares, los índices de enrutamiento evitan tanto un índice central
como inundar la red.
Cada nodo resume, por vecino, cuántos documentos de cada tema se alcanzan a través de él,
y con ese resumen decide a quién reenviar la consulta~\cite{crespo2002routing}; el índice
crece con los vecinos, no con los documentos.
Lo mismo hace en la fase madura cada módulo con su directorio y los de sus conocidos.
En la recuperación de información distribuida hay que elegir, entre muchas bases de
texto, en cuáles buscar~\cite{callan2000distributed}; el directorio de EAM-TMS hace esa
elección con una memoria asociativa.
